\customchapter{INTRODUCCI\'ON}

En los últimos años, el aprendizaje automático distribuido ha cobrado una importancia 
creciente debido al aumento de dispositivos conectados y a la gran cantidad de datos 
generados y almacenados por distintas organizaciones \cite{kairouz2021advances, mahmud2026survey}. 
Esta situación ha impulsado el desarrollo de modelos capaces de aprender a partir de información 
distribuida y de adaptarse a poblaciones con características diversas \cite{DBLP:journals/corr/abs-1908-07873}. 
Sin embargo, en sectores como la salud, la educación y las finanzas, el uso de estos datos está
limitado por normativas de privacidad como HIPAA (\textit{Health Insurance Portability and 
Accountability Act}) y GDPR (\textit{General Data Protection Regulation}), que restringen el 
intercambio de información sensible entre instituciones \cite{noor2026federated}. Frente a este 
desafío, el \textit{Federated Learning} (FL) se ha consolidado como una alternativa viable 
\cite{kairouz2021advances}, ya que permite entrenar modelos de manera local en cada entidad y 
compartir únicamente las actualizaciones de los parámetros con un servidor central encargado de 
combinarlas, evitando así la transferencia directa de los datos originales 
\cite{al2502principles,DBLP:journals/corr/McMahanMRA16}.


A pesar de su elegancia conceptual, el FL enfrenta un desafío fundamental en escenarios reales:
las distribuciones de datos locales difieren significativamente entre clientes, ya sea por 
diferencias en las poblaciones atendidas, en los protocolos de recolección o en los equipos 
utilizados. Esta condición, conocida como \textit{non-IID} (\textit{non-Independent and 
Identically Distributed}), viola el supuesto de homogeneidad subyacente a los análisis de 
convergencia originales del FL. Su consecuencia directa es el fenómeno central que motiva 
esta investigación: el \textit{client drift} \cite{DBLP:journals/corr/abs-1910-06378}. 
Durante el entrenamiento local sobre datos sesgados, el modelo de cada cliente converge hacia 
un óptimo local que puede encontrarse significativamente alejado de la solución global. A 
medida que avanzan las rondas de comunicación, esta divergencia acumulada genera actualizaciones 
que, al promediarse en el servidor, ubican al modelo global en una región de alta pérdida para 
todos los participantes. El resultado es que el progreso de aprendizaje obtenido en cada ronda 
se anula sistemáticamente \cite{seo2025understanding}.


La comunidad investigadora ha respondido a este problema principalmente mediante el desarrollo 
de algoritmos de agregación más robustos, como FedProx \cite{DBLP:journals/corr/abs-1812-06127}, 
SCAFFOLD \cite{DBLP:journals/corr/abs-1910-06378} y sus variantes, que incorporan términos de 
regularización o mecanismos de corrección de varianza para contener la divergencia entre clientes. 
Estos enfoques han logrado mejoras sustanciales en escenarios de heterogeneidad moderada, pero 
su efectividad depende de supuestos sobre el paisaje de pérdida del modelo que no siempre se cumplen 
en la práctica. En particular, la relación entre la \textit{arquitectura} del modelo y la severidad 
con que se manifiesta el \textit{client drift} ha recibido atención considerablemente menor 
\cite{qu2022rethinking, siomos2024architecture}, pese a ser una variable de diseño que los 
profesionales deben definir antes de cualquier elección algorítmica.


Las \textit{Convolutional Neural Networks} (CNNs) explotan filtros espaciales locales; los 
\textit{Vision Transformers} (ViTs) emplean mecanismos de autoatención global; los 
\textit{State-Space Models} (SSMs) modelan dependencias mediante transiciones de estado sobre 
secuencias de parches de la imagen; y los \textit{Vision Graph Neural Networks} (Vision GNNs) 
propagan información mediante mensajes entre regiones de la imagen organizadas en grafos. Estas 
diferencias en los sesgos inductivos implican que cada familia procesa la información de entrada 
de formas radicalmente distintas, lo que debería traducirse en perfiles de drift igualmente 
distintos bajo condiciones de federación heterogénea. Sin embargo, hasta la fecha no existen 
estudios que hayan caracterizado sistemáticamente cómo estas diferencias arquitectónicas se 
manifiestan a nivel de capa a lo largo de las rondas de comunicación, bajo protocolos 
experimentales controlados y comparables entre familias.


Una parte importante de esta brecha es la relacionada con las estrategias de normalización. 
Las CNNs convencionales dependen de la \textit{Batch Normalization} (BN), cuya 
estabilidad requiere estadísticas de mini-batch representativas de la distribución conjunta de 
los datos \cite{DBLP:journals/corr/abs-1803-08494}. En un entorno federado con datos non-IID, 
cada cliente computa estas estadísticas sobre un subconjunto sesgado y reducido, lo que introduce 
una discrepancia entre las estadísticas locales y las globales. Esta discrepancia se propaga 
durante la agregación y puede dominar el comportamiento de convergencia de todo el modelo 
\cite{du2022rethinking}. Por el contrario, las estrategias de normalización independientes del 
tamaño del batch, como \textit{Layer Normalization} (LN) y \textit{Group Normalization} (GN), 
no dependen de las estadísticas entre clientes y se espera que sean inherentemente más estables 
en este contexto \cite{wang2023batch, casella2023experimentingnormalizationlayersfederated}. 
Esta diferencia sugiere que parte de la inestabilidad observada en las CNNs federadas podría 
ser atribuible a la normalización, y no al sesgo inductivo convolucional en sí mismo, pero esta 
hipótesis no ha sido verificada empíricamente a nivel de capa de forma sistemática.


Frente a esta situación, esta investigación propone una caracterización empírica a nivel de capa 
del \textit{client drift} en cuatro familias arquitectónicas de visión: EfficientNet-B1 como 
representante de las CNNs, ViT-Tiny como representante de los Vision Transformers, Vim-Tiny como 
representante de los State-Space Models, y ViG-Tiny (\textit{Vision GNN}) como representante de 
los modelos basados en grafos. Todas las arquitecturas son evaluadas a escala de parámetros 
comparable, bajo condiciones de federación idénticas, en dos conjuntos de datos complementarios: 
Brain Tumor MRI, que representa un escenario realista de imagenología médica con heterogeneidad 
natural entre instituciones, y CIFAR-10, que constituye un punto de referencia estándar en la 
literatura de FL. El grado de heterogeneidad no-IID es controlado mediante particiones de 
Dirichlet con $\alpha \in \{1000.0, 1.0, 0.3, 0.1, 0.03\}$, y se evalúan dos algoritmos de 
agregación: FedAvg y FedProx.


Para caracterizar el drift, se emplean tres métricas complementarias registradas por capa y por 
ronda de comunicación. La divergencia de pesos ($L_2$) mide la distancia entre los parámetros 
locales de cada cliente y el modelo global antes de la agregación, cuantificando la magnitud del 
drift acumulado por tipo de capa. La alineación de gradientes (similitud coseno entre pares de 
clientes) evalúa si las direcciones de actualización de distintos participantes son constructivas 
o destructivas, identificando las capas con mayor interferencia entre clientes durante la 
agregación. El \textit{Centered Kernel Alignment} (CKA) \cite{DBLP:journals/corr/abs-1905-00414} 
compara las representaciones internas aprendidas entre clientes a lo largo del entrenamiento, 
permitiendo determinar si el drift en los pesos se traduce en una divergencia funcional en las 
representaciones, o si los modelos convergen a representaciones similares a pesar de la 
disparidad paramétrica. La combinación de estas tres métricas permite desarrollar un perfil de 
deriva (\textit{drift profile}) para cada arquitectura.


Adicionalmente, se realiza una ablación de normalización sobre EfficientNet, reemplazando su 
BatchNorm por GroupNorm y LayerNorm, con el objetivo de determinar si la inestabilidad federada 
observada en las CNNs es atribuible a la divergencia de estadísticas de lote —un efecto propio 
de la normalización— o al sesgo inductivo convolucional como factor independiente.


Se espera que esta investigación contribuya con: (1) la primera caracterización empírica 
transversal entre familias CNN, ViT, SSM y Vision GNN, con granularidad por capa, sobre cómo 
los sesgos inductivos modulan el \textit{client drift} bajo condiciones non-IID controladas; y 
(2) la incorporación de CKA como métrica de análisis representacional que complementa las medidas 
paramétricas existentes, permitiendo distinguir entre drift en pesos y divergencia funcional en 
las representaciones. Adicionalmente, se aportan: (3) una ablación de normalización que disocia 
la inestabilidad de BatchNorm del sesgo inductivo convolucional; y (4) recomendaciones prácticas, 
fundamentadas en evidencia empírica a nivel de capa, para la selección de arquitecturas y 
estrategias de normalización en escenarios de FL con datos heterogéneos.


\introsection{Descripci\'on del problema}

% -----------------------------------------------------------------------
% CAMBIO: "mecanismos de selección de contenido secuenciales" →
% "transiciones de estado sobre secuencias de parches de la imagen"
% (consistencia con la intro general y más intuitivo para visión).
% "filtrar características" → "privilegiar características"
% ("filtrar" era ambiguo: podía leerse como seleccionar o descartar).
% "practicantes" → "profesionales".
% -----------------------------------------------------------------------
A pesar de su relevancia práctica, la literatura no ha examinado sistemáticamente 
cómo la arquitectura del modelo influye en el \textit{client drift} durante el entrenamiento 
federado con datos heterogéneos \cite{pieri2023handlingdataheterogeneityarchitectural, li2023fedtp}. 
La investigación existente ha abordado este problema casi exclusivamente desde la perspectiva 
del algoritmo de agregación \cite{siomos2024architecture}, asumiendo implícitamente que la 
arquitectura es un factor neutral. Sin embargo, la manera en que cada familia arquitectónica 
procesa la información local determina de forma directa cuánto y de qué manera diverge el 
modelo local del modelo global a lo largo de las rondas de comunicación.

Las CNNs procesan la información mediante filtros locales con pesos compartidos y dependen de 
estadísticas de mini-batch para su normalización, lo que las expone al sesgo local de los datos 
y a la incompatibilidad de estadísticas entre clientes. Los Vision Transformers procesan la 
imagen mediante autoatención global, lo que puede derivar en cabezales de atención fuertemente 
especializados hacia las clases locales bajo distribuciones heterogéneas. Los State-Space Models 
modelan dependencias espaciales mediante transiciones de estado sobre secuencias de parches de 
la imagen; bajo distribuciones non-IID, dichos mecanismos podrían aprender a privilegiar 
características representativas de las clases locales, acumulando sesgo progresivamente a lo 
largo de la secuencia de parches. Los Vision GNNs construyen grafos dinámicos sobre regiones 
de la imagen y propagan información mediante mensajes entre nodos; bajo distribuciones 
heterogéneas, la topología del grafo local podría divergir sistemáticamente entre clientes, 
amplificando el \textit{drift} a través de las capas de propagación. Estos mecanismos 
distintos sugieren que el \textit{drift} varía entre familias no solo en 
magnitud, sino también en su distribución interna por capas y representaciones aprendidas.

A pesar de estas diferencias, no existe un estudio que compare sistemáticamente cómo estas cuatro 
familias experimentan el \textit{client drift}, cómo la 
divergencia se distribuye entre los distintos tipos de capas dentro de cada arquitectura, y 
si la divergencia paramétrica se traduce en una divergencia representacional medible mediante 
CKA. Esta ausencia limita la capacidad de los profesionales de FL para tomar decisiones 
informadas sobre la arquitectura más apropiada para un escenario dado, y restringe el diseño 
de estrategias de agregación que aprovechen la estructura específica del modelo.


\introsection{Justificaci\'on}

La presente investigación responde a una necesidad concreta en el campo del aprendizaje 
federado aplicado a visión por computadora: comprender cómo la elección de arquitectura 
afecta la estabilidad del entrenamiento distribuido bajo condiciones de heterogeneidad 
estadística. Desde una perspectiva científica, contribuye con evidencia empírica que permite 
distinguir entre el efecto del sesgo inductivo y el efecto de la estrategia de normalización 
sobre el \textit{drift}, dos factores que la literatura ha tratado de manera conjunta sin 
separarlos experimentalmente. La incorporación de CKA como métrica complementaria representa 
además un aporte metodológico que enriquece el marco de análisis disponible para estudiar la 
evolución de representaciones internas en sistemas federados.

Desde una perspectiva aplicada, los resultados de esta investigación tienen implicaciones 
directas para cualquier organización que deba implementar FL en entornos con datos heterogéneos 
(como centros de salud, sistemas educativos distribuidos o redes de sensores industriales). 
En estos contextos, seleccionar una arquitectura de visión sin evidencia sobre su comportamiento 
bajo heterogeneidad estadística implica un riesgo de diseño no trivial. Los perfiles de 
\textit{drift} por familia arquitectónica, combinados con los resultados de la ablación de 
normalización, proporcionarán recomendaciones concretas para esas decisiones de diseño.


\introsection{Objetivos}


\introsubsection{Objetivo general}

Desarrollar y comparar perfiles de \textit{drift} a nivel de capa para las familias CNN, 
ViT, SSM y Vision GNN bajo condiciones de federación non-IID controladas, con el fin de 
identificar cómo los sesgos inductivos de cada familia modulan la divergencia de parámetros, 
la alineación de gradientes y la similitud de representaciones internas durante el 
entrenamiento federado.


\introsubsection{Objetivos espec\'ificos}

\begin{enumerate}
    \item \sloppy Medir la divergencia de pesos ($L_2$) por tipo de capa para EfficientNet-B1, 
    ViT-Tiny, Vim-Tiny y ViG-Tiny bajo particiones Dirichlet con 
    $\alpha \in \{1000.0, 1.0, 0.3, 0.1, 0.03\}$, empleando FedAvg y FedProx sobre los 
    conjuntos de datos Brain Tumor MRI y CIFAR-10.

    \item Evaluar la alineación de gradientes entre pares de clientes mediante similitud 
    coseno por capa, identificando las familias y tipos de capa más propensos a interferencia 
    destructiva durante la agregación.

    \item Analizar la evolución de las representaciones internas de cada arquitectura mediante 
    CKA entre pares de clientes a lo largo de las rondas de comunicación, determinando si el 
    \textit{drift} paramétrico implica necesariamente una divergencia funcional en las 
    representaciones aprendidas.

    \item Aislar la contribución de la estrategia de normalización al \textit{drift} 
    observado en EfficientNet mediante una variante experimental que sustituye BatchNorm 
    por GroupNorm y LayerNorm, manteniendo fija la arquitectura base y la configuración 
    de federación.

    \item Formular recomendaciones prácticas para la selección de arquitecturas y estrategias 
    de normalización en escenarios de FL con datos heterogéneos, fundamentadas en los perfiles 
    de \textit{drift} empíricos obtenidos.
\end{enumerate}